\documentclass[11pt]{article}

\pdfinfoomitdate=1
\pdftrailerid{}
\pdfsuppressptexinfo=15

\usepackage[a4paper,top=24mm,bottom=25mm,left=27mm,right=27mm]{geometry}
\usepackage[T1]{fontenc}
\usepackage[utf8]{inputenc}
\usepackage[protrusion=true,expansion=false]{microtype}
\usepackage{amsmath,amssymb}
\usepackage{booktabs}
\usepackage{array}
\usepackage{graphicx}
\usepackage{fancyhdr}
\usepackage[hidelinks]{hyperref}

\hypersetup{
  pdftitle={Cross-Market Coherence between Bitcoin and Nasdaq Futures},
  pdfauthor={Robert Mazurczak},
  pdfsubject={Completed development study, protocol v1.0.0}
}

\setlength{\parindent}{1.25em}
\setlength{\parskip}{0.15em}
\linespread{1.03}
\setlength{\tabcolsep}{5pt}
\renewcommand{\arraystretch}{1.14}

\pagestyle{fancy}
\fancyhf{}
\fancyhead[L]{\footnotesize BTC--NQ cross-market coherence}
\fancyhead[R]{\footnotesize Completed development study}
\fancyfoot[C]{\footnotesize\thepage}
\renewcommand{\headrulewidth}{0.35pt}
\setlength{\headheight}{14pt}
\fancypagestyle{plain}{
  \fancyhf{}
  \fancyhead[L]{\footnotesize BTC--NQ cross-market coherence}
  \fancyhead[R]{\footnotesize Completed development study}
  \fancyfoot[C]{\footnotesize\thepage}
  \renewcommand{\headrulewidth}{0.35pt}
}

\newcommand{\BTC}{\mathrm{BTC}}
\newcommand{\NQ}{\mathrm{NQ}}
\newcommand{\E}{\mathbb{E}}
\newcommand{\I}{\mathbb{I}}
\newcommand{\Mzero}{M_{0}}
\newcommand{\Mone}{M_{1}}
\newcommand{\bps}{\mathrm{bps}}
\newcommand{\statusline}[1]{\noindent\small #1}
\newcolumntype{L}[1]{>{\raggedright\arraybackslash}p{#1}}

\begin{document}

\pagenumbering{Roman}
\begin{titlepage}
  \thispagestyle{empty}
  \vspace*{1.5cm}
  {\raggedright\Huge\bfseries
    Cross-Market Coherence between Bitcoin and Nasdaq Futures\par}
  \vspace{0.35cm}
  {\raggedright\Large
    Magnitude-conditioned state persistence without incremental return predictability\par}
  \vspace{0.85cm}
  \noindent{\large Robert Mazurczak}\par
  \vspace{0.12cm}
  \statusline{September 2026}
  \vspace{0.65cm}
  \hrule
  \vspace{0.5cm}

  \noindent{\large\bfseries Abstract}\par
  \vspace{0.25cm}
  {\small\noindent
  This study asks whether contemporaneous Bitcoin and Nasdaq futures candle bodies
  define persistent cross-market states and whether relative move magnitude adds
  economically useful information. One-minute BTCUSDT and stitched NQ futures bars are
  synchronized in UTC, filtered with DST-aware exchange sessions, and aggregated to a
  five-minute primary representation. Directional coherence is measured over fixed
  15-, 30-, and 60-minute clocks. Body magnitude is normalized by a causal same-slot
  median absolute deviation estimated from 63 prior eligible sessions. A purged,
  expanding walk-forward design compares a coherence-only logistic state model with a
  magnitude-conditioned extension and separately compares NQ-only and cross-market
  Ridge return forecasts. Magnitude reduces five-minute out-of-fold Brier loss by
  $6.99\times10^{-4}$, with a complete-session bootstrap 95\% interval of
  $[-1.120,-0.297]\times10^{-3}$. However, the cross-market return model increases mean
  squared error by $0.1924\,\bps^2$ relative to NQ-only, with interval
  $[0.0728,0.3222]$. The state result repeats overnight. The evidence therefore supports
  a narrow conditional-state association, not incremental NQ return predictability or
  tradable alpha.\par}

  \vspace{0.45cm}
  \noindent\small\textbf{Keywords:} Bitcoin; Nasdaq futures; cross-market dependence;
  directional coherence; walk-forward validation; calibration; data leakage.\par
  \vspace{0.2cm}
  \noindent\textbf{JEL classification:} C53; G14; G17.\par
  \vfill
  \noindent\small\textbf{Research status.} Completed development study, protocol
  v1.0.0. The record was locked after results for reproducibility; it is not a
  preregistration and no pristine final holdout was opened. All reported predictions are
  out of fold. Raw proprietary NQ data are not distributed.
\end{titlepage}
\pagenumbering{arabic}

\section{Introduction}

Bitcoin and Nasdaq futures often print candle bodies in the same direction while their
volatility-adjusted magnitudes differ. That observation is contemporaneous: it does not
by itself imply that BTC leads NQ, that NQ leads BTC, or that the weaker market must catch
up. The empirical question is instead whether synchronized direction defines a persistent
state, whether the scale and imbalance of the joint move help predict that state's next
transition, and whether any such information survives as an executable future NQ return.

The distinction matters. A model can improve the probability forecast for the next
same-direction state without improving a conditional-mean return forecast. Likewise, a
statistically stable state transition can be too small or too indirect to support a
trading rule. This paper therefore separates three estimands: a descriptive response
curve, a binary state-persistence forecast, and an incremental NQ return forecast.

The contribution is methodological rather than a claim of alpha. The analysis enforces
next-open timing, causal normalization, fold-local preprocessing, nested regularization,
complete-session purging, block-bootstrap inference, and a mandatory overnight negative
control. Trial ledgers and compact aggregate outputs are retained so that unfavorable
comparisons cannot be silently discarded. This follows the broader concern that repeated
selection can make backtest performance materially overstate evidence
\cite{bailey2014deflated,bailey2017pbo,lopezdeprado2018}.

\section{Research questions and estimands}

Let $O^i_t$ and $C^i_t$ denote the open and close of asset $i\in\{\BTC,\NQ\}$ during bar
$t$. The signed candle-body return is
\begin{equation}
  b^i_t = \log\!\left(\frac{C^i_t}{O^i_t}\right).
  \label{eq:body}
\end{equation}
Wicks and high--low ranges are excluded from the primary definition. Per-bar directional
agreement is
\begin{equation}
  d_t = \operatorname{sign}(b^{\BTC}_t)\operatorname{sign}(b^{\NQ}_t)
  \in\{-1,0,+1\}.
  \label{eq:agreement}
\end{equation}
The study addresses two primary questions. First, conditional on $d_t=+1$, do magnitude
features improve the forecast of $d_{t+1}=+1$ beyond recent directional coherence?
Second, do BTC and coherence variables improve the forecast of the executable next
five-minute signed NQ return beyond an otherwise comparable NQ-only information set?
Longer horizons and the one-minute representation are prespecified secondary analyses.

\section{Data and sample construction}

The available sample combines Binance BTCUSDT perpetual one-minute OHLCV with a local
stitched NQ front-month series. The common verified range is 7 March 2024 through
12 May 2026. Timestamps are stored in UTC and mapped to \texttt{America/New\_York} for
session logic. The cash-session sample follows the official 09:30--16:00 New York equity
session, including early closes. The overnight control covers 18:00 on the prior evening
through 09:30, excluding the CME maintenance interval, weekends, holidays, gaps, and
contract-transition sessions.

\begin{table}[!htbp]
  \centering
  \caption{Data provenance and eligibility constraints}
  \label{tab:data}
  \small
  \begin{tabular}{L{0.20\textwidth}L{0.28\textwidth}L{0.39\textwidth}}
    \toprule
    Input & Available coverage & Principal constraint \\
    \midrule
    BTCUSDT perpetual & 2024-03-07 to 2026-05-12 & Public Binance archive; post-ETP sample only \\
    NQ stitched front month & Common range through 2026-05-12 & Partial contract provenance; historical roll rule cannot be reconstructed \\
    Cash sessions & 458 complete sessions before warm-up & Entire incomplete or transition sessions excluded \\
    Overnight sessions & 551 complete sessions before warm-up & Separate negative control; cannot replace cash primary \\
    \bottomrule
  \end{tabular}
  \par\vspace{2pt}
  \parbox{0.90\textwidth}{\footnotesize\textit{Notes:} Prices are never forward-filled.
  Five-minute bars require all constituent one-minute observations. A 63-session causal
  warm-up is applied after eligibility.}
\end{table}

The data begin after US spot Bitcoin ETP trading and were inspected in earlier project
iterations. Consequently, the motivating post-ETP integration narrative is not tested as
a causal before/after design, and this sample cannot provide a pristine final holdout.
These restrictions are part of the inferential result rather than implementation details.

\section{Feature construction}

\subsection{Directional coherence}

For clock-time window $K\in\{15,30,60\}$ minutes, coherence is the unweighted average of
complete agreement observations,
\begin{equation}
  \mathcal{C}_t(K)=\frac{1}{n_K}\sum_{j=0}^{n_K-1}d_{t-j},
  \qquad
  \mathbf{C}_t=\bigl[\mathcal{C}_t(15),\mathcal{C}_t(30),\mathcal{C}_t(60)\bigr].
  \label{eq:coherence}
\end{equation}
The windows remain continuous on $[-1,1]$ and are used jointly. No historically best
window or coherence threshold is selected. The same clocks correspond to 3/6/12 bars in
the five-minute representation and 15/30/60 bars in the one-minute robustness analysis.

\subsection{Causal body scale and magnitude coordinates}

Let $\mathcal{H}^{i,r}_t$ contain the previous 63 eligible body returns for asset $i$,
resolution $r$, and the same DST-aware session slot as $t$. The current session is
excluded. Define
\begin{align}
  \widetilde b^{i,r}_{t-1}
    &= \operatorname{median}_{x\in\mathcal{H}^{i,r}_t} x, \\
  \widehat\sigma^{i,r}_{t-1}
    &= 1.4826\operatorname{median}_{x\in\mathcal{H}^{i,r}_t}
       \left|x-\widetilde b^{i,r}_{t-1}\right|, \\
  m^i_t &= \frac{|b^i_t|}{\widehat\sigma^{i,r}_{t-1}}.
  \label{eq:scale}
\end{align}
An observation is ineligible when fewer than 63 valid references exist or the scale is
zero or non-finite. There is no epsilon floor or future-data fallback.

Joint intensity and relative magnitude balance are
\begin{equation}
  J_t=\sqrt{m^{\BTC}_t m^{\NQ}_t},
  \qquad
  B_t=\frac{m^{\BTC}_t-m^{\NQ}_t}{m^{\BTC}_t+m^{\NQ}_t}.
  \label{eq:magnitude}
\end{equation}
$J_t$ is symmetric and becomes large only when both markets move strongly relative to
their own histories. $B_t\in[-1,1]$ separates relative strength from common event scale;
positive values indicate a relatively larger BTC body.

\begin{table}[!htbp]
  \centering
  \caption{Frozen parameter provenance}
  \label{tab:parameters}
  \small
  \begin{tabular}{L{0.27\textwidth}L{0.18\textwidth}L{0.44\textwidth}}
    \toprule
    Parameter & Value & Role and status \\
    \midrule
    Primary bars & 5 minutes & Tractable primary representation; not claimed optimal \\
    Robustness bars & 1 minute & Mandatory timing check with the same economic endpoint \\
    Coherence clocks & 15/30/60 minutes & Multiscale state vector; no outcome-based best-window selection \\
    Historical scale & 63 sessions & Causal same-slot MAD; current session excluded \\
    Economic horizon & 5 minutes & Next-open to horizon-open at both resolutions \\
    Purge & 1 complete session & Exceeds the target horizon and separates adjacent folds \\
    \bottomrule
  \end{tabular}
\end{table}

\section{Prediction design and inference}

\subsection{Causal timing and outcomes}

Features use information through the event-bar close $\tau_t$. Economic measurement
begins at the first tradable open after signal formation and ends at a future open:
\begin{equation}
  R^{\NQ}_{t,h}=\log\!\left(\frac{O^{\NQ}_{\tau_t+h}}{O^{\NQ}_{\tau_t}}\right),
  \qquad
  Y_{t,h}=S_t R^{\NQ}_{t,h},
  \label{eq:return}
\end{equation}
where $S_t$ is the common detected direction and $h=5$ minutes is primary. Thus positive
$Y_{t,h}$ denotes continuation and negative values denote reversal. Horizons 15, 30, and
60 minutes form a secondary response curve; none can replace the five-minute endpoint.

\subsection{Nested model comparisons}

Conditional on $d_t=+1$, the state models are logistic regressions,
\begin{align}
  \Mzero &: \Pr(d_{t+1}=+1\mid\mathbf{C}_t,S_t), \\
  \Mone &: \Pr(d_{t+1}=+1\mid\mathbf{C}_t,S_t,J_t,B_t,|B_t|).
  \label{eq:state-models}
\end{align}
The primary state score is Brier loss; log loss and fixed-width reliability bins are
secondary. The return comparison uses matched Ridge models. The NQ-only information set
contains current NQ direction and magnitude, NQ momentum at 5/15/30/60 minutes, intraday
Fourier terms, and weekday controls. The cross-market model adds BTC momentum and
magnitude, the coherence vector, $J_t$, $B_t$, and $|B_t|$.

Both comparisons use expanding outer folds with 130 initial feature-ready sessions,
one complete purged session, and 12 consecutive validation blocks of 22 sessions. All
preprocessing and regularization selection occur within folds. Logistic
$C\in\{0.01,0.1,1,10\}$ and Ridge
$\alpha\in\{0.01,0.1,1,10,100,1000\}$ are selected from three purged inner expanding
blocks. Every scored prediction is out of fold.

Paired uncertainty resamples complete validation sessions 5,000 times. For model loss
$L$, the reported incremental comparison is
\begin{equation}
  \Delta L = L_{\mathrm{cross}}-L_{\mathrm{baseline}}.
  \label{eq:loss-difference}
\end{equation}
Negative values favor the expanded model. This convention is used for both the Brier and
return-loss comparisons.

\section{Empirical results}

\subsection{Descriptive response}

\begin{figure}[!htbp]
  \centering
  \includegraphics[width=0.92\textwidth]{report/figures/response_curve.png}
  \caption{Signed NQ open-to-open response conditional on current BTC--NQ body-direction
  agreement. Intervals are 95\% complete-session bootstrap intervals. The five-minute
  cash-session endpoint is primary; longer horizons and overnight estimates are
  secondary.}
  \label{fig:response}
\end{figure}

The cash-session five-minute descriptive estimate is $0.012\,\bps$ across 19,955 events,
with an interval containing zero. The positive longer-horizon point estimates are much
less precise and do not alter the registered primary conclusion. The unconditional
response therefore supplies no trading premise before model comparison.

\subsection{State persistence}

The five-minute state sample contains 13,724 out-of-fold events from 264 validation
sessions. The observed next-bar agreement rate is 68.25\%. Table~\ref{tab:state-results}
shows that adding magnitude reduces both Brier and log loss.

\begin{table}[!htbp]
  \centering
  \caption{Five-minute out-of-fold state-prediction results}
  \label{tab:state-results}
  \small
  \begin{tabular}{lrrr}
    \toprule
    Model & Brier loss & Log loss & ECE \\
    \midrule
    Fold training-rate benchmark & 0.217250 & 0.626229 & 0.009849 \\
    $\Mzero$: coherence & 0.216736 & 0.625038 & 0.010755 \\
    $\Mone$: coherence + magnitude & \textbf{0.216037} & \textbf{0.623370} & \textbf{0.009218} \\
    \midrule
    $\Mone-\Mzero$ & $-0.000699$ & $-0.001668$ & $-0.001537$ \\
    95\% block interval & $[-0.001120,-0.000297]$ & $[-0.002668,-0.000703]$ & --- \\
    \bottomrule
  \end{tabular}
  \par\vspace{2pt}
  \parbox{0.88\textwidth}{\footnotesize\textit{Notes:} ECE denotes expected calibration
  error. Lower values are better. The Brier reduction is approximately 0.32\% relative
  to $\Mzero$.}
\end{table}

$\Mone$ improves Brier loss in 10 of 12 outer folds. Standardized coefficient signs are
stable: mean coefficients are approximately $+0.10$ for joint intensity, $-0.06$ for
signed balance, and $+0.03$ for absolute imbalance. These are conditional predictive
associations, not structural or causal effects.

\subsection{Incremental return value}

\begin{figure}[!htbp]
  \centering
  \includegraphics[width=0.94\textwidth]{report/figures/primary_results.png}
  \caption{Primary five-minute paired loss differences. Points are out-of-fold means;
  horizontal bars are 95\% complete-session bootstrap intervals. Negative values favor
  the expanded cross-market model. The panels use different loss units.}
  \label{fig:primary}
\end{figure}

The return sample contains 11,364 events from the same 264 validation sessions. Neither
Ridge model beats the fold training-mean benchmark, and the expanded information set is
worse than NQ-only under both registered losses.

\begin{table}[!htbp]
  \centering
  \caption{Five-minute signed NQ return forecast}
  \label{tab:return-results}
  \small
  \begin{tabular}{lrrrr}
    \toprule
    Model & MSE ($\bps^2$) & RMSE ($\bps$) & MAE ($\bps$) & OOS $R^2$ \\
    \midrule
    Fold training mean & \textbf{134.831} & \textbf{11.612} & \textbf{7.080} & 0.0000 \\
    NQ-only Ridge & 135.101 & 11.623 & 7.092 & -0.0020 \\
    Cross-market Ridge & 135.294 & 11.632 & 7.099 & -0.0034 \\
    \midrule
    Cross-market $-$ NQ-only & $+0.1924$ & --- & $+0.0073$ & --- \\
    95\% block interval & $[+0.0728,+0.3222]$ & --- & $[+0.0010,+0.0135]$ & --- \\
    \bottomrule
  \end{tabular}
  \par\vspace{2pt}
  \parbox{0.88\textwidth}{\footnotesize\textit{Notes:} Lower forecast loss is better.
  The positive paired differences favor NQ-only. OOS $R^2$ is relative to the fold-local
  training-mean forecast.}
\end{table}

The registered incremental-return criterion is therefore not met. Strong shrinkage is
itself informative: $\alpha=1000$ is selected in 9 of 12 NQ-only folds and 8 of 12
cross-market folds. The nested process generally prefers predictions close to the local
mean instead of large feature-driven variation.

\section{Robustness and negative control}

At one-minute resolution, $\Mone-\Mzero$ Brier loss is $-0.000297$ with 95\% interval
$[-0.000418,\allowbreak -0.000173]$ across 67,483 events. The five-minute economic endpoint remains
fixed, however, and the cross-market minus NQ-only return MSE difference is
$+0.0031\,\bps^2$ with interval $[-0.1065,\allowbreak +0.0962]$. The finer representation therefore
supports the narrow state result but does not rescue the return hypothesis.

The overnight state sample contains 36,230 observations. Its $\Mone-\Mzero$ Brier
difference is $-0.000259$ with interval $[-0.000440,-0.000081]$. Cash-session specificity
is therefore not supported. The overnight return comparison remains economically
unhelpful: cross-market minus NQ-only MSE is $+0.0106\,\bps^2$ with interval
$[-0.0283,+0.0502]$, while MAE worsens and both Ridge models trail the fold mean.

\begin{figure}[!htbp]
  \centering
  \includegraphics[width=0.86\textwidth]{report/figures/state_calibration.png}
  \caption{Reliability diagrams for the five-minute cash-session state model and the
  overnight negative control. Empty probability bins are omitted from the lines but
  retained in the committed calibration tables.}
  \label{fig:calibration}
\end{figure}

\section{Limitations and interpretation}

The evidence is deliberately narrower than a typical trading-strategy claim.
\begin{itemize}
  \item The sample has no pre-ETP history, so the motivating market-integration narrative
  is neither identified nor tested causally.
  \item No pristine holdout exists. Intervals condition on the fitted out-of-fold
  predictions and do not quantify uncertainty from repeating the entire research process
  on an independent history.
  \item NQ is a partial stitched series. The complete historical roll-selection rule and
  contract-level lineage cannot be reconstructed.
  \item Missing sessions are systematic; notably few Thursdays survive. External
  validity cannot assume missingness at random.
  \item NQ price behavior is not MNQ execution evidence. No transaction-cost, liquidity,
  slippage, position-sizing, or deployment claim is made.
\end{itemize}

The state result is statistically stable within the development sample but practically
small. Its repetition overnight weakens a cash-session-specific interpretation. Most
importantly, the return test demonstrates that improved state classification need not
translate into improved conditional-mean return prediction. Opening an entry/exit search
after that result would increase researcher degrees of freedom without an economic basis.

\section{Conclusion}

Volatility-adjusted joint magnitude adds a small, repeatable amount of information about
whether synchronized BTC and NQ candle-body directions persist into the next bar. It does
not improve the registered five-minute NQ return forecast over an NQ-only baseline. The
appropriate conclusion is conditional state dependence without demonstrated alpha.

This negative economic result is useful: it rejects the simplest tradability route while
preserving a precisely defined state variable for future work. A confirmatory extension
would require complete contract-level data, a genuinely untouched period, and a new
protocol fixed before any nonlinear model, event filter, or coherence-decay policy is
evaluated.

\clearpage
\section*{Reproducibility and credit}
\thispagestyle{fancy}

Code, configurations, trial ledgers, aggregate results, figures, and SHA-256 manifests
are versioned in the public repository. Raw proprietary NQ rows remain local. The
research question originated in earlier collaborative coursework with Alex Samuseva;
the present repository is a clean redesign and does not treat legacy notebook results as
evidence. No license has been selected pending a joint authorship decision.

\begin{thebibliography}{9}
\thispagestyle{fancy}

\bibitem{bailey2014deflated}
Bailey, D. H., and L{\'o}pez de Prado, M. (2014).
The Deflated Sharpe Ratio: Correcting for Selection Bias, Backtest Overfitting, and
Non-Normality. \emph{Journal of Portfolio Management}, 40(5), 94--107.

\bibitem{bailey2017pbo}
Bailey, D. H., Borwein, J. M., L{\'o}pez de Prado, M., and Zhu, Q. J. (2017).
The Probability of Backtest Overfitting. \emph{Journal of Computational Finance},
20(4), 39--69.

\bibitem{brier1950}
Brier, G. W. (1950).
Verification of Forecasts Expressed in Terms of Probability.
\emph{Monthly Weather Review}, 78(1), 1--3.

\bibitem{efron1993}
Efron, B., and Tibshirani, R. J. (1993).
\emph{An Introduction to the Bootstrap}. Chapman \& Hall/CRC.

\bibitem{lopezdeprado2018}
L{\'o}pez de Prado, M. (2018).
\emph{Advances in Financial Machine Learning}. Wiley.

\end{thebibliography}

\end{document}
